\documentclass[../report.tex]{subfiles}
\begin{document}

\section{Phần mở rộng đã làm}

Ngoài yêu cầu cơ bản, nhóm làm thêm các phần sau:

\begin{itemize}
  \item Các lệnh tùy chọn của đề: LB, LH, LBU, LHU, SB, SH, BLT, BGE, BLTU, BGEU, cùng LUI và AUIPC.
  \item GPIO và UART theo kiểu memory-mapped (\sig{IOController}, mục nâng cao của đề). Chương trình
        \file{t11\_io.s} đọc 8 công tắc DIP, ghi LED, nhận ký tự từ Keyboard, đổi sang chữ hoa và in ra TTY.
  \item Forwarding đủ hai nguồn EX/MEM và MEM/WB, bypass WB$\to$ID, và tín hiệu \sig{UseRs1}, \sig{UseRs2} để không
        stall khi lệnh không thật sự đọc thanh ghi.
  \item Bộ đếm \sig{cycles}, \sig{instret} trong mạch và Pipeline Monitor hiển thị lệnh ở từng tầng.
  \item Toolchain: assembler, golden model, bộ chạy test tự động không cần giao diện, bộ sinh chương trình ngẫu nhiên,
        và công cụ \file{tools/compare.py} in giản đồ pipeline của bất kỳ chương trình nào.
  \item Bản NH và chế độ assembler \texttt{--nohazard}, để đo phần đóng góp của phần cứng xử lý hazard.
\end{itemize}

\section{Hạn chế}

\begin{itemize}
  \item Truy cập bộ nhớ lệch (ví dụ \texttt{lw} ở địa chỉ không chia hết cho 4) không được hỗ trợ; kết quả không
        xác định. Golden model báo lỗi trong trường hợp này.
  \item Không có exception và CSR. ECALL, EBREAK được dùng làm lệnh dừng; FENCE thực thi như NOP.
  \item Thanh ghi và PC reset về 0; phần mềm phải tự đặt \texttt{sp}.
  \item Vùng I/O chỉ giải mã \sig{Addr[31:28]} và \sig{Addr[4:2]}, nên mọi địa chỉ \texttt{0xFxxx\_xxxx} khác là bí
        danh của 8 thanh ghi I/O.
  \item UART không có ngắt, phần mềm phải hỏi vòng thanh ghi trạng thái; Keyboard chỉ nhận ký tự ASCII 7 bit.
  \item SB, SH dùng đọc, sửa, ghi thay cho byte enable của RAM. Trong Logisim cách này đúng vì RAM đọc tổ hợp; trên
        phần cứng thật nó làm đường store dài hơn.
  \item Rẽ nhánh quyết định ở EX nên mỗi lần nhảy mất 2 chu kỳ. Trên 12 chương trình test, flush chiếm
        $2\times\resTotBR$ chu kỳ, nhiều hơn tổng \resTotLU{} chu kỳ của load-use.
  \item Logisim không mô phỏng trễ cổng, nên mọi so sánh thời gian trong chương~\ref{ch:result} dựa trên giả định
        trễ từng tầng, không phải số đo.
\end{itemize}

\section{Hướng phát triển}\label{sec:future}

Các mục nâng cao chưa làm gồm dự đoán nhánh, cache, nhân chia, ngắt và số thực dấu phẩy động. Bảng~\ref{tab:future}
xếp chúng theo lợi ích có thể đo bằng bộ test hiện có.

\begin{table}[H]
\centering\small
\caption{Hướng phát triển và cách đo lợi ích}\label{tab:future}
\begin{tabularx}{\textwidth}{@{}>{\raggedright\arraybackslash}p{3.6cm}>{\raggedright\arraybackslash}p{4.8cm}X@{}}
\toprule
Mở rộng & Cách làm & Lợi ích có thể đo \\
\midrule
Quyết định rẽ nhánh ở ID & Đưa \sig{BranchUnit} và bộ cộng địa chỉ đích lên ID, thêm forwarding vào ID & Phạt mỗi lần nhảy giảm từ 2 xuống 1 chu kỳ. Trên \file{t09\_recursion.s} tiết kiệm tối đa \resNineBR{} chu kỳ; phải trừ đi các stall mới khi lệnh rẽ nhánh phụ thuộc lệnh ngay trước \\
Dự đoán nhánh 2 bit và BTB & Bảng nhỏ đánh chỉ số bằng \sig{PC[5:2]}, lưu địa chỉ đích và bộ đếm bão hòa 2 bit & Giảm flush ở vòng lặp (\file{t08\_sort.s}, \file{t09\_recursion.s}); đo bằng số lần \sig{E\_Redirect} \\
Phần mở rộng M (MUL, DIV) & Thêm Multiplier, Divider của Logisim vào ALU, chọn theo funct7 = 0000001 & \file{t09\_recursion.s} đang nhân bằng vòng lặp dịch và cộng; thay bằng một lệnh \texttt{mul} \\
Cache & Cache ánh xạ trực tiếp trước \sig{DMEM}, thêm tín hiệu stall khi miss & Cần mô hình độ trễ bộ nhớ chính mới đo được \\
Ngắt và CSR & Thêm \texttt{mtvec}, \texttt{mepc}, \texttt{mcause} và các lệnh CSR & UART nhận ký tự không cần hỏi vòng \\
\bottomrule
\end{tabularx}
\end{table}

\end{document}
